\section{Configuration model}

\subsection{Profiles}

A profile directory contains one JSON file per enabled subsystem. The aggregated prototype iterates over the registered builders and includes only files that exist:

\begin{lstlisting}
config/profiles/default/
  global.json
  psu.json
  chiller.json
\end{lstlisting}

The full simulated profile is in \file{config/profiles/prototype}. Raspberry hardware uses \file{config/profiles/raspberry}.

The command:

\begin{lstlisting}[style=shell]
bdx-prototype-ioc --config-dir config/profiles/default
\end{lstlisting}

builds all supported subsystem files found in that directory. Duplicate PV names are rejected. All included subsystem configurations must resolve to the same server interface list.

\subsection{Common server settings}

A typical JSON file contains:

\begin{lstlisting}[style=json]
{
  "server": {
    "interfaces": ["0.0.0.0"],
    "poll_interval": 1.0,
    "log_pv_names": false
  }
}
\end{lstlisting}

\code{BDX_EPICS_INTERFACE} overrides \code{server.interfaces} at runtime. This is the preferred mechanism on a host with more than one network interface.

The configured polling interval must be positive. In the aggregated prototype, all managed IOC groups share one \code{RuntimeSettings} object; the global IOC can change the effective update period at runtime within configured bounds.

\subsection{Device entries}

Hardware-specific addresses, ports, timeouts, channel lists, optional features, and limits belong in JSON. IOC classes must not hard-code those deployment values.

Example PSU device entry:

\begin{lstlisting}[style=json]
{
  "name": "LV1",
  "prefix": "BDX:PSU:LV1:",
  "mode": "hardware",
  "driver": "cpx400dp",
  "host": "172.22.50.20",
  "port": 9221,
  "timeout": 3.0,
  "channels": [1, 2],
  "software_limits": {
    "minimum_voltage": 0.0,
    "maximum_voltage": 60.0,
    "minimum_current_limit": 0.0,
    "maximum_current_limit": 20.0,
    "maximum_power": 420.0
  }
}
\end{lstlisting}

Device limits may be overridden per channel through \code{channel_limits}. Validate a new profile with the simulated or disconnected startup path before connecting it to hardware.

\subsection{Runtime environment}

Machine-specific values are not committed in \file{config/runtime.env}. Create it from \file{config/runtime.env.example} and set the operational main IOC address:

\begin{lstlisting}[style=shell]
BDX_MAIN_HOST=172.22.50.2
\end{lstlisting}

Do not store credentials, Wi-Fi secrets, private keys, downloaded third-party runtimes, logs, databases, or archive data in the repository.

\section{PV contract}

\subsection{Naming convention}

The default hierarchy is:

\begin{lstlisting}
BDX:<SUBSYSTEM>:<DEVICE>:<CHANNEL>:<FIELD>
\end{lstlisting}

Examples:

\begin{lstlisting}
BDX:PSU:LV1:CH1:VOLTAGE_REQUEST
BDX:PSU:LV1:CH1:VOLTAGE_RBV
BDX:CHILLER:CHILLER1:SETPOINT_REQUEST
BDX:ENV:TEMP:T00:VALUE
BDX:DAQ:CRATE01:READY
BDX:GLOBAL:INTERLOCK_ACTIVE
\end{lstlisting}

Use these suffixes consistently:

\begin{tabularx}{\textwidth}{@{}l X@{}}
\toprule
Suffix & Meaning \\
\midrule
\code{_RBV} & Read-only value obtained from device or IOC state \\
\code{_SET} & Direct writable setpoint or compatibility control \\
\code{_REQUEST} & Operator-staged value not yet applied to hardware \\
\code{_CMD} & Momentary command, normally self-resetting \\
\code{_STATUS} & Textual outcome or state \\
\code{_OK} & Boolean status suited to indicators and logic \\
\bottomrule
\end{tabularx}

A successful write means the IOC accepted the command path; it does not prove that the physical device reached the requested state. Update readbacks only from driver-reported state.

\subsection{Managed IOC health PVs}

Every class derived from \code{ManagedIOC} exposes:

\begin{lstlisting}
HEARTBEAT
IOC_STATE
COMM_STATUS
COMM_OK
LAST_UPDATE
ERROR_CODE
ERROR_MESSAGE
SIMULATION
CLEAR_ERROR_CMD
\end{lstlisting}

At startup, \code{IOC_STATE} changes to \code{RUNNING}. The heartbeat loop performs a driver communication check, polls subsystem state, updates health PVs, and sleeps for the shared runtime period. Communication failures set \code{COMM_OK=false} and preserve an actionable error message.

The baseline communication states are \code{STARTING}, \code{OK}, \code{DISCONNECTED}, \code{TIMEOUT}, \code{INVALID_DATA}, \code{DEVICE_ERROR}, \code{DISABLED}, and \code{SIMULATION}. New drivers should map failures consistently rather than inventing display-specific semantics.
